% GENERATED FILE. Edit canonical lessons, then run npm run generate:books.
% canonical-source-hash: fnv1a64:65cd2326da82a3d3
% canonical-lessons: TE-C544-kridakarudu, TE-C544-pratyarthi, TE-C544-talapadu, TE-C544-neggu, TE-C544-parugu

\chapter{Sport}
\label{ch:te-sport}
\hlchaptermodality{544}

\begin{chapteropening}
\textbf{By the end of this chapter:} \emph{I can talk about sportsmen, rivals, matches, wins and running.}

Complete the last lesson of chapter 544: I can talk about sportsmen, rivals, matches, wins and running.
\end{chapteropening}

\section[\texorpdfstring{\te{క్రీడాకారుడు} (krīḍākāruḍu)}{క్రీడాకారుడు (krīḍākāruḍu)}]{\te{క్రీడాకారుడు} (krīḍākāruḍu) — a sportsman; an athlete}

\label{lesson:TE-C544-kridakarudu}



\begin{quote}
Before the new one: say the Telugu for to savour, then the Telugu for a stick dance.
\end{quote}

\subsection*{You'll want to know: \te{క్రీడాకారుడు}}
\textbf{\te{క్రీడాకారుడు}} — \emph{krīḍākāruḍu} — \textquotedblleft{}a sportsman; an athlete\textquotedblright{}.

\begin{quote}
\te{ఆ} \te{క్రీడాకారుడు} \te{దేశానికి} \te{పతకం} \te{తెచ్చాడు}.
\end{quote}

\begin{quote}
\emph{ā krīḍākāruḍu dēśāniki patakaṁ teccāḍu.}
\end{quote}

\begin{quote}
That sportsman won a medal for the country.
\end{quote}

\te{క్రీడ}, sport, + -\te{కారుడు}, one who does it.

\begin{cousinweb}
\textbf{\te{క్రీడాకారుడు}} (\emph{krīḍākāruḍu}) is built on Sanskrit \emph{krīḍā}, \textquotedblleft{}sport, play, pastime\textquotedblright{} (Monier-Williams), with -\te{కారుడు}, \textquotedblleft{}one who does\textquotedblright{}.
\end{cousinweb}

\begin{grammarlens}[title={What you've built}]
One more word for describing what you lived and felt.
\end{grammarlens}

\subsection*{Your turn}
\begin{itemize}
\raggedright
  \item \emph{Say it:} \emph{krīḍākāruḍu}
  \item \emph{Say it:} \emph{krīḍākāruḍu}, once more
  \item \emph{Say it:} \emph{ā krīḍākāruḍu dēśāniki patakaṁ teccāḍu.} — That sportsman won a medal for the country.
  \item \emph{Recall:} say the Telugu for to savour, then the Telugu for a stick dance, then say \emph{krīḍākāruḍu} again
\end{itemize}

\begin{tcolorbox}[breakable,colback=teal!4,colframe=teal!35!black,title={Before you move on}]
What does \te{క్రీడాకారుడు} mean? (\textquotedblleft{}A sportsman; an athlete\textquotedblright{}.) Say it once more.
\end{tcolorbox}
\section[\texorpdfstring{\te{ప్రత్యర్థి} (pratyarthi)}{ప్రత్యర్థి (pratyarthi)}]{\te{ప్రత్యర్థి} (pratyarthi) — an opponent; a rival}

\label{lesson:TE-C544-pratyarthi}



\begin{quote}
Before the new one: say the Telugu for a stick dance, then the Telugu for a sportsman.
\end{quote}

\subsection*{You'll want to know: \te{ప్రత్యర్థి}}
\textbf{\te{ప్రత్యర్థి}} — \emph{pratyarthi} — \textquotedblleft{}an opponent; a rival\textquotedblright{}.

\begin{quote}
\te{మా} \te{ప్రత్యర్థి} \te{చాలా} \te{బలమైన} \te{జట్టు}.
\end{quote}

\begin{quote}
\emph{mā pratyarthi cālā balamaina jaṭṭu.}
\end{quote}

\begin{quote}
Our opponent is a very strong team.
\end{quote}

In sport, politics or business.

\begin{cousinweb}
\textbf{\te{ప్రత్యర్థి}} (\emph{pratyarthi}) is borrowed from Sanskrit \emph{pratyarthin}, \textquotedblleft{}an adversary, an opponent, a rival\textquotedblright{} (Monier-Williams).
\end{cousinweb}

\begin{grammarlens}[title={What you've built}]
One more word for describing what you lived and felt.
\end{grammarlens}

\subsection*{Your turn}
\begin{itemize}
\raggedright
  \item \emph{Say it:} \emph{pratyarthi}
  \item \emph{Say it:} \emph{pratyarthi}, once more
  \item \emph{Say it:} \emph{mā pratyarthi cālā balamaina jaṭṭu.} — Our opponent is a very strong team.
  \item \emph{Recall:} say the Telugu for a stick dance, then the Telugu for a sportsman, then say \emph{pratyarthi} again
\end{itemize}

\begin{tcolorbox}[breakable,colback=teal!4,colframe=teal!35!black,title={Before you move on}]
What does \te{ప్రత్యర్థి} mean? (\textquotedblleft{}An opponent; a rival\textquotedblright{}.) Say it once more.
\end{tcolorbox}
\section[\texorpdfstring{\te{తలపడు} (talapaḍu)}{తలపడు (talapaḍu)}]{\te{తలపడు} (talapaḍu) — to clash; to take on}

\label{lesson:TE-C544-talapadu}



\begin{quote}
Before the new one: say the Telugu for a sportsman, then the Telugu for an opponent.
\end{quote}

\subsection*{You'll want to know: \te{తలపడు}}
\textbf{\te{తలపడు}} — \emph{talapaḍu} — \textquotedblleft{}to clash; to take on\textquotedblright{}.

\begin{quote}
\te{రేపు} \te{రెండు} \te{జట్లు} \te{ఫైనల్లో} \te{తలపడతాయి}.
\end{quote}

\begin{quote}
\emph{rēpu reṇḍu jaṭlu phainallō talapaḍatāyi.}
\end{quote}

\begin{quote}
Tomorrow the two teams clash in the final.
\end{quote}

Past (I …): \te{తలపడ్డాను}, \emph{talapaḍḍānu}.

-\te{తో} \te{తలపడు}, to take someone on.

\begin{cousinweb}
\textbf{\te{తలపడు}} (\emph{talapaḍu}) is a native compound: \emph{tala}, \textquotedblleft{}head\textquotedblright{} (DEDR 3103), and the native verb \emph{paḍu}, \textquotedblleft{}to feel, to suffer\textquotedblright{} (DEDR 3853): to meet head to head.
\end{cousinweb}

\begin{grammarlens}[title={What you've built}]
One more word for describing what you lived and felt.
\end{grammarlens}

\subsection*{Your turn}
\begin{itemize}
\raggedright
  \item \emph{Say it:} \emph{talapaḍu}
  \item \emph{Say it:} \emph{talapaḍu}, once more
  \item \emph{Say it:} \emph{rēpu reṇḍu jaṭlu phainallō talapaḍatāyi.} — Tomorrow the two teams clash in the final.
  \item \emph{Recall:} say the Telugu for a sportsman, then the Telugu for an opponent, then say \emph{talapaḍu} again
\end{itemize}

\begin{tcolorbox}[breakable,colback=teal!4,colframe=teal!35!black,title={Before you move on}]
What does \te{తలపడు} mean? (\textquotedblleft{}To clash; to take on\textquotedblright{}.) Say it once more.
\end{tcolorbox}
\section[\texorpdfstring{\te{నెగ్గు} (neggu)}{నెగ్గు (neggu)}]{\te{నెగ్గు} (neggu) — to prevail; to win through}

\label{lesson:TE-C544-neggu}



\begin{quote}
Before the new one: say the Telugu for an opponent, then the Telugu for to clash.
\end{quote}

\subsection*{You'll want to know: \te{నెగ్గు}}
\textbf{\te{నెగ్గు}} — \emph{neggu} — \textquotedblleft{}to prevail; to win through\textquotedblright{}.

\begin{quote}
\te{చివరి} \te{బంతికి} \te{మా} \te{జట్టు} \te{నెగ్గింది}.
\end{quote}

\begin{quote}
\emph{civari bantiki mā jaṭṭu neggindi.}
\end{quote}

\begin{quote}
Our team won on the last ball.
\end{quote}

Past (I …): \te{నెగ్గాను}, \emph{neggānu}.

Also in arguments and elections; \te{గెలుచు} is the everyday word.

\begin{cousinweb}
\textbf{\te{నెగ్గు}} (\emph{neggu}) is native Dravidian: DEDR 3661 lists Telugu \emph{neggu}, \textquotedblleft{}to be successful, to succeed, to thrive\textquotedblright{}.
\end{cousinweb}

\begin{grammarlens}[title={What you've built}]
One more word for describing what you lived and felt.
\end{grammarlens}

\subsection*{Your turn}
\begin{itemize}
\raggedright
  \item \emph{Say it:} \emph{neggu}
  \item \emph{Say it:} \emph{neggu}, once more
  \item \emph{Say it:} \emph{civari bantiki mā jaṭṭu neggindi.} — Our team won on the last ball.
  \item \emph{Recall:} say the Telugu for an opponent, then the Telugu for to clash, then say \emph{neggu} again
\end{itemize}

\begin{tcolorbox}[breakable,colback=teal!4,colframe=teal!35!black,title={Before you move on}]
What does \te{నెగ్గు} mean? (\textquotedblleft{}To prevail; to win through\textquotedblright{}.) Say it once more.
\end{tcolorbox}
\section[\texorpdfstring{\te{పరుగు} (parugu)}{పరుగు (parugu)}]{\te{పరుగు} (parugu) — running; a sprint}

\label{lesson:TE-C544-parugu}



\begin{quote}
Before the new one: say the Telugu for to clash, then the Telugu for to prevail.
\end{quote}

\subsection*{You'll want to know: \te{పరుగు}}
\textbf{\te{పరుగు}} — \emph{parugu} — \textquotedblleft{}running; a sprint\textquotedblright{}.

\begin{quote}
\te{ఉదయం} \te{పరుగు} \te{ఆరోగ్యానికి} \te{మంచిది}.
\end{quote}

\begin{quote}
\emph{udayaṁ parugu ārōgyāniki mañcidi.}
\end{quote}

\begin{quote}
A morning run is good for your health.
\end{quote}

The noun; \te{పరిగెత్తు}, to run, is the verb.

\begin{cousinweb}
\textbf{\te{పరుగు}} (\emph{parugu}) is native Dravidian: DEDR 3963 lists Telugu \emph{parugu}, \textquotedblleft{}running, a run\textquotedblright{}, in the article of Tamil \emph{pari}, \textquotedblleft{}to run, to go fast\textquotedblright{}.
\end{cousinweb}

\begin{grammarlens}[title={What you've built}]
One more word for describing what you lived and felt.
\end{grammarlens}

\subsection*{Your turn}
\begin{itemize}
\raggedright
  \item \emph{Say it:} \emph{parugu}
  \item \emph{Say it:} \emph{parugu}, once more
  \item \emph{Say it:} \emph{udayaṁ parugu ārōgyāniki mañcidi.} — A morning run is good for your health.
  \item \emph{Recall:} say the Telugu for to clash, then the Telugu for to prevail, then say \emph{parugu} again
\end{itemize}

\begin{tcolorbox}[breakable,colback=teal!4,colframe=teal!35!black,title={Before you move on}]
What does \te{పరుగు} mean? (\textquotedblleft{}Running; a sprint\textquotedblright{}.) Say it once more.
\end{tcolorbox}
